\documentclass[a4paper,11pt]{article}
\usepackage{exam-style-341}
\examinehead{341 农业综合知识三 · 网络技术与应用}{模拟试卷（一）参考答案}
\begin{document}

\answercover{341 农业综合知识三 · 网络技术与应用}{模拟试卷（一）参考答案}

\answerpart{一、单项选择题}
\setcounter{question}{0}
\q C。\quad \q B。\quad \q B。\quad \q B。\quad \q C。\quad \q C。\quad \q A。\quad \q B。\quad \q B。\quad \q B。
\begin{analysis}
第 1 题：时延带宽积 \(=\) 传播时延 \(\times\) 带宽，表示链路中「正在传播的比特数」，又称以比特为单位的链路长度。
第 2 题：数据链路层在相邻结点间的链路上传输以帧为单位的数据，通过 CRC 检错、确认与重传实现无差错传输。
第 3 题：个人区域网（PAN，约 10 m）\(\rightarrow\) 局域网（LAN）\(\rightarrow\) 城域网（MAN）\(\rightarrow\) 广域网（WAN）。
第 4 题：B 类地址第一个字节为 \(128\sim191\)。A 类为 \(1\sim126\)，C 类为 \(192\sim223\)，\(226\) 属 D 类（多播）。
第 5 题：光纤靠光信号传输，不受电磁干扰，频带宽、损耗小、保密性好；无线电波易受干扰，双绞线和同轴电缆带宽较小。
第 6 题：数字信号在普通电话线上远距离传输必须先经调制解调器（MODEM）变换为模拟信号，C 错误。
第 7 题：ARP 把 IP 地址解析为 MAC 地址；RARP 反向解析；ICMP 用于差错与询问报文；IGMP 用于多播组管理。
第 8 题：用户代理把邮件交给发送方邮件服务器、以及邮件服务器之间传送邮件都用 SMTP；接收邮件用 POP3 或 IMAP，POP3 使用 110 端口，IMAP 在服务器上管理邮件，故只有 B 正确。
第 9 题：对等层之间传送的数据单元就是该层的协议数据单元 PDU。
第 10 题：IP 提供无连接的、尽最大努力交付的服务，不保证可靠、不保证按序、无端到端流量控制，可靠性由运输层（如 TCP）负责。
\end{analysis}

\answerpart{二、填空题}
\setcounter{question}{0}
\q 语义；同步（时序）。【两个空顺序可互换，写「同步（时序）；语义」同样给分。】\quad
\q 广域网。\quad
\q 波分复用（WDM）；码分复用（CDM，又称码分多址 CDMA）。【顺序可互换；写「空分复用」「统计时分复用」等教材所列其他复用技术亦给分。】\quad
\q \(1\sim126\)（二进制首位为 0）；\(255.0.0.0\)。【第一空写 \(1\sim126\)、\(1\) 至 \(126\)、128 种网络号等均可；写 \(1\sim127\) 或 \(0\sim126\) 不给分。】\quad
\q RTT（往返时间）。
\begin{analysis}
第 1 题：语法规定数据与控制信息的格式，语义规定需要发出何种控制信息以及完成何种动作，同步（时序）规定事件实现顺序的详细说明。
第 4 题：A 类地址共 32 位，第 1 位固定为 0，故第一个字节最大为 \(01111111_2=127\)；其中 \(127\) 保留作环回地址，可用网络号范围为 \(1\sim126\)；默认掩码为 \(11111111.00000000.00000000.00000000=255.0.0.0\)。
第 5 题：TCP 每发送一个报文段就记录其发送时间，收到确认后计算 RTT 的加权平均值（Karn 算法等），再据此计算出 RTO，故 RTT 是确定 RTO 的主要依据。
\end{analysis}

\answerpart{三、名词解释}
\setcounter{question}{0}
\q \textbf{基带信号}：来自信源（如计算机、话筒）的、未经调制的原始电信号，其频谱从零频附近开始，包含较低的频率成分（含直流分量）。基带信号可以直接在电缆上传输，称为基带传输；也可以先经调制变换为带通信号，在无线信道或电话线上进行频带传输。

\q \textbf{ARP（地址解析协议，Address Resolution Protocol）}：把 IP 地址解析为硬件（MAC）地址的网络层协议。主机发送数据前先在本机 ARP 高速缓存中查找目的 IP 对应的 MAC 地址；找不到时以广播方式发送 ARP 请求，目标主机以单播方式返回含自身 MAC 地址的 ARP 响应，并把结果存入缓存。ARP 只能在同一局域网（同一广播域）内使用，路由器不转发 ARP 请求。

\q \textbf{FTP（文件传输协议，File Transfer Protocol）}：互联网上使用最广泛的文件传送协议，基于 TCP 提供可靠传输。它采用客户/服务器方式，工作时使用两条并行的连接：控制连接使用 21 端口，在整个会话期间保持；数据连接使用 20 端口（主动方式），仅在传送文件时建立，传送完毕即关闭。FTP 的控制信息与文件数据分开传送（带外控制），服务器进程可同时为多个客户服务。

\q \textbf{RTT（往返时间，Round-Trip Time）}：从发送方发送数据开始，到发送方收到来自接收方的确认（接收方收到数据后立即回送确认）总共经历的时间，是衡量网络时延与性能的重要指标。RTT 由链路的传播时延、结点的排队与处理时延及确认时间等组成，它随网络负载变化。TCP 用 RTT 的测量值按加权平均等方式估算超时重传时间 RTO；ping 命令也用 RTT 表示连通性与时延大小。

\answerpart{四、简答与计算题}
\setcounter{question}{0}

\q \textbf{（1）CSMA/CD 工作过程（载波监听多点接入/碰撞检测，3 分）}
\begin{enumerate}[leftmargin=2em, itemsep=2pt, topsep=3pt]
  \item \textbf{先听后发（载波监听）}：每个站点在发送数据前先监听总线，若总线忙则继续监听，直到总线变为空闲；
  \item \textbf{边听边发（碰撞检测）}：发送过程中继续监听，把自己发出的信号与总线上收到的信号比较，以判断是否发生了碰撞；
  \item \textbf{冲突停发}：一旦检测到碰撞，立即停止发送，并发送若干比特的人为干扰信号（Jamming Signal），使所有站点都能检测到碰撞；
  \item \textbf{随机重发（二进制指数退避）}：停止发送后按二进制指数退避算法等待一段随机时间，再重新尝试发送，重传次数越多，随机等待时间的均值越大。
\end{enumerate}

\textbf{（2）争用期（碰撞窗口）与最短帧长（2 分）}
\begin{itemize}[leftmargin=2em, itemsep=2pt, topsep=3pt]
  \item 争用期是发送方发出数据后，在最远距离的站点处仍可能发生碰撞的最长时间，其值为端到端传播时延的两倍（\(2\tau\)）。如果发送方在争用期内没有检测到碰撞，才能确信本次发送不会发生碰撞；因此以太网端到端往返时间 \(2\tau\) 称为争用期（碰撞窗口），10 Mbit/s 以太网规定争用期为 \(51.2\ \mu s\)。
  \item 为使发送站在发送完毕之前就能检测到碰撞，数据帧的发送时延必须不小于争用期，即最短帧长 \(=\) 争用期 \(\times\) 数据率 \(=51.2\ \mu s\times 10\ \mathrm{Mbit/s}=512\ \mathrm{bit}=64\) 字节。凡长度小于 64 字节的帧都是因碰撞而异常中止的无效帧，直接丢弃，从而使 CSMA/CD 能够正常工作。
\end{itemize}
\begin{analysis}
本题为简答题，答出「先听后发、边听边发、冲突停发、随机重发」四个步骤即得 3 分；说明争用期为 \(2\tau\)、最短帧长 \(=\) 争用期 \(\times\) 数据率 \(=64\) 字节及其作用得 2 分。
\end{analysis}

\q \textbf{解：}

\textbf{（1）子网位数与子网掩码（2 分）}

网络地址 \(168.195.0.0\) 第一个字节为 \(168\in[128,191]\)，是 B 类地址，其标准（默认）子网掩码为 \(255.255.0.0\)，即网络号为 16 位、主机号为 16 位。

设借用 \(n\) 位主机号作为子网号，按 \(2^{n}\ge 27\) 计算：
\[
2^{4}=16<27,\qquad 2^{5}=32\ge 27
\]
故至少需要借用 \(n=5\) 位作为子网号，这样最多可划分 \(2^{5}=32\) 个子网（其中全 0 和全 1 的子网号一般不用，实际可用 30 个），满足划分 27 个子网的要求。

新的子网掩码为「默认掩码的网络号部分 \(16\) 位 \(+\) 借用的子网号 \(5\) 位 \(=21\) 个连续的 1」，即：
\[
11111111.11111111.11111000.00000000=255.255.248.0
\]

\textbf{（2）每个子网可容纳的主机数（1.5 分）}

用去 21 位后，主机号还剩 \(32-21=11\) 位，因此每个子网的地址总数为 \(2^{11}=2048\) 个。其中主机号全 0 的地址是该子网的网络地址，主机号全 1 的地址是该子网的广播地址，二者都不能分配给主机，所以每个子网最多可容纳的主机数为：
\[
2^{11}-2=2048-2=2046\ \text{台}
\]
题目中 \(2^{11}=2048\) 是子网内的地址总数，不是可分配的主机数，必须减去网络地址和广播地址这 2 个特殊地址。

\textbf{（3）子网掩码的二进制形式与第 1 个可用子网（1.5 分）}

子网掩码 \(255.255.248.0\) 的二进制展开为：
\[
\underbrace{11111111}_{255}.\underbrace{11111111}_{255}.\underbrace{11111000}_{248}.\underbrace{00000000}_{0}
\]
第三个字节中，前 5 位（\(11111\)）是子网号，后 3 位（\(000\)）与第 4 字节一起构成 11 位主机号。子网号部分按 \(00001\)、\(00010\)、\(\cdots\) 顺序取值（\(00000\) 为子网号全 0 的子网，不使用）：

\begin{center}
\begin{tabular}{cccc}
\toprule
子网号（二进制） & 第三个字节 & 子网网络地址 & 可分配给主机的地址范围 \\
\midrule
\(00000\) & 0 & \(168.195.0.0\) & 不使用（子网号全 0） \\
\(00001\) & 8 & \(168.195.8.0\) & \(168.195.8.1\sim168.195.15.254\) \\
\(00010\) & 16 & \(168.195.16.0\) & \(168.195.16.1\sim168.195.23.254\) \\
\(00011\) & 24 & \(168.195.24.0\) & \(168.195.24.1\sim168.195.31.254\) \\
\(\cdots\) & \(\cdots\) & \(\cdots\) & \(\cdots\) \\
\(11110\) & 240 & \(168.195.240.0\) & \(168.195.240.1\sim168.195.247.254\) \\
\(11111\) & 248 & \(168.195.248.0\) & 不使用（子网号全 1） \\
\bottomrule
\end{tabular}
\end{center}

所以第 1 个可用子网的网络地址是 \(\mathbf{168.195.8.0}\)（掩码 \(255.255.248.0\)），广播地址是 \(168.195.15.255\)，可分配给主机的地址范围是 \(\mathbf{168.195.8.1\sim168.195.15.254}\)，共 2046 个可用地址。

\textbf{验算}：子网块大小 \(=256-248=8\)，即第三个字节以 8 为步长递增，\(8\times 31=248\)，共 32 个子网；子网 \(168.195.8.0/21\) 的有效地址为 \(256\times 8-2=2048-2=2046\) 个，与（2）一致。

\begin{analysis}
本题评分要点：①由 \(2^n\ge 27\) 说明 \(n=5\)（写 \(2^4<27\le 2^5\) 给满分），并正确写出掩码 \(255.255.248.0\)；②正确写出 \(2^{11}-2=2046\)，且说明要减去网络地址与广播地址；\par ③正确写出掩码的二进制展开 \(11111111.11111111.11111000.00000000\)\allowbreak，并给出第 1 个可用子网 \(168.195.8.0\) 及地址范围 \(168.195.8.1\sim168.195.15.254\)。常见错误：只借 4 位（\(2^4=16<27\)）；忘记减 2，把 \(2048\) 当可用主机数；把子网号全 0 的子网 \(168.195.0.0\) 当作第 1 个可用子网。
\end{analysis}

\end{document}
